% TOP_PROBLEM: {"id": 3000011, "title": "Chromatic number of t-perfect graphs", "category_id": 2, "category": {"id": 2, "name": "combinatorics", "display_name": "Combinatorics", "description": "Counting problems, graph theory, discrete structures.", "slug": "combinatorics", "order_index": 2, "created_at": "2026-07-31T15:26:25.671Z"}, "status": "open", "problem_number": "AMR-029-0011", "set_id": 13}
% FIRST_POSTED: 2026-10-01
% ENTRY_KIND: statement_only
% UnsolvedMath ID 3000011; source catalogue code AMR-029-0011
% !TeX program = lualatex
% Definition and sources only; this file is not an LLM solution attempt.
\documentclass[11pt,a4paper]{article}
\usepackage[margin=25mm]{geometry}
\usepackage{fontspec}
\setmainfont{lmroman10-regular.otf}[BoldFont=lmroman10-bold.otf,
 ItalicFont=lmroman10-italic.otf,BoldItalicFont=lmroman10-bolditalic.otf]
\usepackage{amsmath,amssymb,mathtools}
\usepackage{unicode-math}
\setmathfont{latinmodern-math.otf}
\AtBeginDocument{\let\setminus\smallsetminus}
\usepackage{xurl,hyperref}
\hypersetup{colorlinks=true,urlcolor=blue}
\setcounter{secnumdepth}{0}
\setlength{\parindent}{0pt}
\setlength{\parskip}{5pt}
\setlength{\emergencystretch}{3em}
\begin{document}
\section*{Chromatic number of t-perfect graphs}\label{3000011}
{\small UnsolvedMath ID 3000011 · AMR-029-0011 · Combinatorics · AMR Open Problem Lists}
\subsection{Definitions and mathematical statement}
Let $G=(V,E)$ be a finite simple graph and write $x(U)=\sum_{v\in U}x_v$.
A stable set is a vertex set containing no edge, and
$\operatorname{STAB}(G)$ is the convex hull of its characteristic vectors.
Define the odd-cycle relaxation
\[
\begin{split}
 Q(G)=\{x\in[0,1]^V:\;&x_u+x_v\le1\quad(uv\in E),\\
 &x(V(C))\le (|V(C)|-1)/2\quad\text{for every odd simple cycle }C\}.
\end{split}
\]
The graph is called $t$-perfect if $Q(G)=\operatorname{STAB}(G)$.
The coordinate upper bounds make the definition apply also to isolated
vertices, whose presence does not affect the colouring question.

Is every $t$-perfect graph properly colourable with at most four colours?
Explicitly, does there always exist a function
\[
                  f:V\longrightarrow\{1,2,3,4\},\qquad
                  f(u)\ne f(v)\quad(uv\in E)?
\]
Equivalently, can $V$ be partitioned into at most four stable sets?
No planarity, forbidden-claw, or triangle-free assumption is included.
The number four is absolute and must not grow with the number of vertices.
The defining polyhedral equality concerns fractional stable-set constraints;
converting it into a bounded partition of all vertices is the additional
requirement. A bound on the fractional chromatic number or a larger fixed
colour bound does not by itself answer this four-colour conjecture.
\subsection{Short English statement}
Must every graph whose stable-set polytope is described by edge
and odd-cycle inequalities admit a proper colouring with four colours?
\subsection{Sources}
\begin{itemize}
\item[S1] ULAM.AI and the UnsolvedMath Contributors, supplied problems.json, record 3000011 (AMR-029-0011), AMR Open Problem Lists. Catalogue identity and source transcription; CC BY 4.0.
\url{https://huggingface.co/datasets/ulamai/UnsolvedMath}.
\item[S2] Egres Open - List of open problems, Open-problem page: Chromatic number of t-perfect graphs. Original source indexed by AMR-029-0011.
\url{https://oldlemon.cs.elte.hu/egres/open/List_of_open_problems}.
\item[S3] EGRES Open, Chromatic number of t-perfect graphs. Individual source page and explanatory remarks.
\url{https://lemon.cs.elte.hu/egres/open/Chromatic_number_of_t-perfect_graphs}.
\end{itemize}
\end{document}
